%%%%%%%%%%  scan idx 119  |  folio 106  %%%%%%%%%%
\SGAPrintPage{106}%
so \(\Ho^n(F'_{\leq n})\) (resp.\ \(\Ho^n(F''_{\leq n})\)) is of finite type, so again, by (2.10.b)), \(F'_{\leq n}\) (resp.\ \(F''_{\leq n}\)) is \(n\)-pseudo-coherent. Applying (2.5) to the exact sequence
\[
  0\longrightarrow F'_{>n}\longrightarrow F'\longrightarrow F'_{\leq n}\longrightarrow 0,
\]
(resp.\dots), we deduce that \(F'\) (resp.\ \(F''\)) is \(n\)-pseudo-coherent, QED.

\sgaStableTarget{sga6.e01.prop.2-13}\textbf{Proposition 2.13.} \uline{Let} \((\underline{S}',\underline{C}',\underline{C}'_0)\) \uline{be a triple satisfying the same hypotheses} (2.0) \uline{as} \((\underline{S},\underline{C},\underline{C}_0)\)\uline{. Let} \(f:\underline{S}'\to \underline{S}\) \uline{be a morphism of sites} ([G] I 1.4)\uline{, corresponding to an underlying functor} \(f^{-1}:\underline{S}\to \underline{S}'\)\uline{, and let}
\(u:\Dpx^*(\underline{C})\longrightarrow \Dpx^*(\underline{C}')\) \((*=+,-,\varnothing)\)
\uline{be a cartesian functor over} \(f\)\uline{,} i.e.\ \uline{a cartesian} \(\underline{S}\)\uline{-functor} \(\Dpx^*(\underline{C}')\times_{\underline{S}'}\underline{S}\longleftarrow\Dpx^*(\underline{C})\)\uline{. Assume that} \(u\) \uline{is exact on each fibre and that there exist} \(a,b\in\mathbb Z\) \uline{such that the following conditions hold for every} \(S\in\Ob \underline{S}\):

(i) \uline{when} \(E\in\Ob\underline{C}_{0S}\)\uline{,} \(u(E)\) \uline{is} \(a\)\uline{-}\(\underline{C}'_0\)\uline{-pseudo-coherent (resp.} \(\underline{C}'_0\)\uline{-pseudo-coherent)}

(ii) \uline{when} \(F\in\Ob\Dpx^*(\underline{C}_S)\) \uline{is acyclic in degree} \(>0\)\uline{,} \(u(F)\) \uline{is acyclic in degree} \(>b\)\uline{.}

\uline{Then, when} \(F\in\Ob\Dpx^*(\underline{C}_S)\) \uline{is} \(m\)\uline{-pseudo-coherent rel.\ to} \(\underline{C}_0\) \uline{and acyclic in degree} \(>n\)\uline{,} \(u(F)\in\Ob\Dpx^*(\underline{C}'_{S'})\) \uline{(where} \(S'=f^{-1}(S)\)\uline{) is} \(\sup(m+b,n+a)\)\uline{-pseudo-coherent (rel.\ to} \(\underline{C}'_0\)\uline{) (resp.} \((m+b)\)\uline{-pseudo-coherent) and acyclic in degree} \(>n+b\)\uline{.}

\uline{In particular, in the case of the ``resp.'' alternative, when} \(F\in\Ob\Dpx^*(\underline{C}_S)\) \uline{is pseudo-coherent} (rel.\ to \(\underline{C}_0\)), \(u(F)\) \uline{is pseudo-coherent (rel.\ to} \(\underline{C}'_0\)\uline{).}

\uline{Proof.} We can restrict ourselves to the case without the ``resp.'' alternative. Let us first show that, for \(F\in\Ob\Dpx^*(\underline{C}_S)\) strictly perfect (2.1) with \(F^i=0\) for \(i>n\), \(u(F)\) is \((n+a)\)-\(\underline{C}'_0\)-pseudo-coherent. We argue by induction on the number \(N\) of indices \(i\) such that \(F^i\neq0\). For \(N=1\),
%%%%%%%%%%  scan idx 120  |  folio 107  %%%%%%%%%%
\SGAPrintPage{107}%
we have \(F=F^p[-p]\) for some \(p\leq n\), so \(u(F)\) is \((p+a)\)-\(\underline{C}'_0\)-pseudo-coherent, hence a fortiori \((n+a)\)-pseudo-coherent. For \(N>1\), there is an exact sequence
\[
  0\longrightarrow F'\longrightarrow F\longrightarrow F''\longrightarrow 0,
\]
where \(F'\) and \(F''\) are strictly perfect of length \(\leq N-1\) and satisfy \(F'^i=F''^i=0\) for \(i>n\). Applying (2.5.b)) to the distinguished triangle
\[
\begin{tikzcd}
& u(F'') \arrow[dl,"+1"'] & \\
u(F') \arrow[rr] && u(F) \arrow[ul] & ,
\end{tikzcd}
\]
we deduce from the induction hypothesis that \(u(F)\) is \((n+a)\)-pseudo-coherent.

Let us now pass to the general case: \(F\) \(m\)-\(\underline{C}_0\)-pseudo-coherent and acyclic in degree \(>n\). After localizing, we can, by (2.10.a)), suppose that \(F\) is strictly \(m\)-pseudo-coherent and satisfies \(F^i=0\) for \(i>n\). There is then an exact sequence
\[
  0\longrightarrow F'\longrightarrow F\longrightarrow F''\longrightarrow 0,
\]
where \(F''\) is zero in degree \(\geq m\), and \(F'\) strictly perfect and zero in degree \(>n\). In the distinguished triangle
\[
\begin{tikzcd}
& u(F'') \arrow[dl,"+1"'] & \\
u(F') \arrow[rr] && u(F) \arrow[ul] & ,
\end{tikzcd}
\]
\(u(F'')\) is acyclic in degree \(\geq m+b\), hence \((m+b)\)-pseudo-coherent, and \(u(F')\) is \((n+a)\)-pseudo-coherent by what we have seen. We conclude by (2.5.b)).

\sgaStableTarget{sga6.e01.cor.2-14}\textbf{Corollary 2.14.} \uline{Let} \(\underline{C}_1\) \uline{be a strictly full additive sub-}\(\underline{S}\)\uline{-category of} \(\underline{C}\) \uline{such that} \(\underline{C}_0\subset\underline{C}_1\subset\underline{C}\)\uline{. Assume that every object of} \(\underline{C}_1\) \uline{is} \(\underline{C}_0\)\uline{-pseudo-coherent.}

%%%%%%%%%%  scan idx 121  |  folio 108  %%%%%%%%%%
\SGAPrintPage{108}%
\uline{Then}:

(i) \(\underline{C}_1\) \uline{is locally quasi-liftable in} \(\underline{C}\) (1.2) \uline{and locally quasi-stable under kernels of epimorphisms} (1.2)\uline{.}

(ii) \uline{Let} \(S\in\Ob \underline{S}\) \uline{and} \(F\in\Ob\Dpx(\underline{C}_S)\)\uline{. For} \(F\) \uline{to be} \(n\)\uline{-pseudo-coherent (resp.\ pseudo-coherent) relative to} \(\underline{C}_1\)\uline{, it is necessary and sufficient that it be so relative to} \(\underline{C}_0\)\uline{.}

\uline{Proof.} (i) Let \(E\to F\) be an epimorphism of \(\underline{C}_S\), with \(F\in\Ob\underline{C}_{1S}\). Since \(F\) is \(\underline{C}_0\)-pseudo-coherent (hence a fortiori of \(\underline{C}_0\)-finite type), there exists (after localizing) an epimorphism \(E'\to F\) with \(E'\in\Ob\underline{C}_{0S}\). Using the quasi-liftability of \(\underline{C}_0\) in \(\underline{C}\), we can suppose (after localizing again) that this factors through \(E\), whence the quasi-liftability of \(\underline{C}_1\). On the other hand, when \(E\in\Ob\underline{C}_S\) admits a finite right resolution by objects of \(\underline{C}_1\), \(E\) is \(\underline{C}_0\)-pseudo-coherent by (2.11.a)), hence a fortiori of \(\underline{C}_1\)-finite type.

(ii) It suffices to prove necessity. For this we can apply, as we choose, (2.11.a)) or (2.13) with \(\underline{S}'=\underline{S}\), \(\underline{C}'=\underline{C}\), \(f=\text{Identity}\), \(u=\text{Identity}\), replacing \(\underline{C}'_0\) and \(\underline{C}_0\) respectively by \(\underline{C}_0\) and \(\underline{C}_1\).

\sgaStableTarget{sga6.e01.schol.2-14-1}\textbf{Scholium 2.14.1.} There thus exists a largest strictly full additive sub-\(\underline{S}\)-category of \(\underline{C}\) giving the same notion of \(n\)-pseudo-coherence (resp.\ pseudo-coherence) as \(\underline{C}_0\), namely the category formed by the \(\underline{C}_0\)-pseudo-coherent objects. The latter is locally stable under kernels of epimorphisms.

\sgaStableTarget{sga6.e01.ex.2-15}\textbf{Example 2.15.} Let \((\underline{S},A)\) be a ringed topos and \(\underline{C},\underline{C}_0\) the \(\underline{S}\)-categories considered in (1.3.d)). The fibered category \(\Dpx(\underline{C})\) will often be denoted by \(\Dpx(\underline{S},A)\), or \(\Dpx({}_A \underline{S})\) (to
%%%%%%%%%%  scan idx 122  |  folio 109  %%%%%%%%%%
\SGAPrintPage{109}%
suggest that these are left \(A\)-Modules), or even \(\Dpx(\underline{S})\) when there is no ambiguity concerning \(A\). We shall denote by \(S\) the final object of \(\underline{S}\) and put \(\Dpx(\underline{C}_S)=\Dpx({}_A S)\) (or \(\Dpx(S)\)). An object of \(\Dpx(\underline{S})\) will be called \(\underline{n}\)\uline{-pseudo-coherent} (resp.\ \uline{pseudo-coherent}) when it is \(n\)-\(\underline{C}_0\)-pseudo-coherent (resp.\ \(\underline{C}_0\)-pseudo-coherent) (likewise for the strict notions). The fibered category \(\Dpx(\underline{C})_{n\text{-}\underline{C}_0\text{-}\coh}\) (resp.\ \(\Dpx(\underline{C})_{\underline{C}_0\text{-}\coh}\)) (2.6) will be denoted by \(\Dpx({}_A \underline{S})_{n\text{-}\coh}\) (resp.\ \(\Dpx({}_A \underline{S})_{\coh}\)). Observe that, by (2.14), we obtain the same notion of pseudo-coherence on replacing \(\underline{C}_0\) by the \(\underline{S}\)-category \(\underline{C}_1\) defined by \(U\longmapsto\) (category of left \(A_U\)-Modules locally direct factors of free Modules of finite type). When an \(A\)-Module \(M\) is a direct factor of a free \(A\)-Module of finite type, it does not follow, in general, that \(M\) is locally free of finite type. However, we have:

\sgaStableTarget{sga6.e01.lem.2-15-1}\textbf{Lemma 2.15.1.} \uline{Assume that} \(\underline{S}\) \uline{satisfies one of the following two hypotheses:}

(i) \(\underline{S}\) \uline{``has enough points''} (SGA 4 IV 6.4) \uline{and, for every point} \(s\) \uline{of} \(\underline{S}\)\uline{, the projective} \(A_s\)\uline{-modules} (on the left) \uline{of finite type are free, this last condition being satisfied for example} (BOURBAKI Alg. Comm. chap.\ II \S~3 n\textsuperscript{o}~2 cor.~2 of prop.\ 5) \uline{when} \(A_s/\underline{\operatorname{radical}}(A_s)\) \uline{is a field.}

(ii) \((\underline{S},A)\) \uline{is a ``locally ringed'' topos in the sense of} [1]\uline{, which means that} \(A\) \uline{is commutative and that, for every} \(X\in\Ob \underline{S}\) \uline{and every} \(f\in\Gamma(X,A)\)\uline{, we have} \(X=X_f\cup X_{1-f}\)\uline{, where} \(X_f\) \uline{(resp.} \(X_{1-f}\)\uline{) denotes the largest subobject of} \(X\) \uline{on which} \(f\) \uline{(resp.} \(1-f\)\uline{) is invertible.}

\uline{Then every direct factor of a free} \(A\)\uline{-module (on the left) of finite type is locally free of finite type.}

\uline{Proof.} Suppose (i). We first observe that, for two \(A\)-Modules \(E\) and \(F\), with \(E\) of finite presentation, the canonical arrow
%%%%%%%%%%  scan idx 123  |  folio 110  %%%%%%%%%%
\SGAPrintPage{110}%
\[
  \underline{\Hom}(E,F)_s\longrightarrow\Hom(E_s,F_s)
\]
is an isomorphism for every point \(s\) of \(\underline{S}\). It follows that, when \(E\) and \(F\) are both of finite presentation, and \(E_s\simeq F_s\) at a point \(s\) of \(\underline{S}\), there exists an \(s\)-pointed object \(U\) of \(\underline{S}\)\SGAAuthorFootnote{*}{i.e.\ such that \(s\) factors through \(\underline{S}_{/U}\).} such that \(E|U\simeq F|U\). The conclusion follows at once.

Suppose (ii). Let \(n\in\mathbb N\), and denote by \(\mathbb D^n\) the functor on the category of locally ringed topoi (belonging to a fixed universe\dots) defined by \(T\longmapsto\Gamma(T,\underline{O}_T)^n\). We know [1] that the functor \(\underline{\End}_{\mathbb D}(\mathbb D^n)\) is represented by \(M_n=\Spec\mathbb Z[T_{ij}]\ (1\leq i\leq n,\ 1\leq j\leq n)\), equipped with the endomorphism \(u_n\) of \((\underline{O}_{M_n})^n\) defined by the matrix \((T_{ij})\). Now let \(p:A^n\longrightarrow A^n\) be a projector; let us show that \(p(A^n)\) is locally free of finite type. By what we have just said, there exists a morphism \(s:\underline{S}\longrightarrow M_n\) of locally ringed topoi such that \(s^*(u_n)=p\). We can interpret \(s\) as a point of \(M_n\) with values in \(\Gamma(S,A)\), i.e.\ a matrix of order \(n\) with coefficients in \(\Gamma(S,A)\); since this matrix is a projector, \(s\) factors through the closed subscheme \(i:P_n\to M_n\) corresponding to projectors as \(\underline{S}\xrightarrow{s'}P_n\xrightarrow{i}M_n\). Thus \(p=s'^*i^*(u_n)\), and we are done, because by the assertion under hypothesis (i), the image of \(i^*(u_n)\) is a locally free sheaf of finite type.

\sgaStableTarget{sga6.e01.cor.2-16}\textbf{Corollary 2.16.} \uline{Let} \(f:(\underline{S}',A')\to(\underline{S},A)\) \uline{be a morphism of ringed topoi (defined by a functor} \(f^{-1}:\underline{S}\to \underline{S}'\) \uline{and a morphism of sheaves of rings} \(f^{-1}(A)\to A'\)\uline{). Denote by} \(\Dpx({}_{A'}\underline{S}'_A)\) \uline{the derived category of the category of} \(A'\)\uline{-}\(A\)\uline{-biModules (on the left over} \(A'\)\uline{, on the right over} \(A\)\uline{), and by}
\[
  \otimes_A^{\mathrm L}:\Dpx^-(_{A'}\underline{S}'_A)\times_{\underline{S}'}\Dpx^-(_A \underline{S})
  \longrightarrow \Dpx^-(_{A'}\underline{S}')
\]
